% !TeX root = ../main.tex
\chapter{Limitations and Future Scope}

No production-oriented civic platform is complete without explicit recognition of current boundaries. Jan-Sunwai AI delivers a strong operational baseline, but several limitations remain due to time, infrastructure, and deployment-context constraints [16, 7, 17]. Current limitations were documented to provide a realistic boundary for what the present release can guarantee in production-like environments. These limitations do not invalidate the solution; rather, they identify the most important areas where additional engineering and policy integration can strengthen long-term adoption.

\section{Current Limitations}

The primary limitations of the current system fall into three categories: infrastructure ceilings, model variability, and institutional adoption challenges. As a closed-system deployment, local inference is bound by available VRAM, meaning high-concurrency requests can bottleneck the synchronous AI pipeline and cause unpredictable response times. Classification quality also remains highly dependent on the clarity of the uploaded civic imagery, meaning ambiguous scenes still heavily rely on the manual triage queue. Additionally, while the application layer is robust against common web vulnerabilities, external perimeter security (such as WAFs and zero-trust segmentation) remains entirely dependent on the municipal host environment.

\section{Future Scope}

The future roadmap focuses on evolving the platform from a pilot-grade application to a highly scalable, institutional-grade ecosystem. Priority enhancements include migrating the synchronous AI pipeline to a distributed, asynchronous message broker to handle municipal-scale concurrency without dropping requests. Furthermore, extending the rule engine to incorporate fine-tuned, localized vision models (rather than relying on zero-shot inference) will drastically reduce misclassification rates. From an operational perspective, future versions should integrate automated SLA policy enforcement with built-in escalation paths to external authority APIs, alongside comprehensive disaster recovery runbooks to ensure uninterrupted civic service.

\section{Chapter Summary}

The current system is robust for pilot and structured deployment but remains intentionally conservative in distributed scaling and advanced automation. The future scope presented here provides a realistic pathway from project-grade implementation to institutional-grade civic infrastructure.
